\documentclass[10pt,a4paper]{amsart}
\usepackage[margin=19mm]{geometry}
\usepackage{amsmath,amssymb,mathtools}
\usepackage[scr=rsfs,cal=euler]{mathalfa}
\usepackage{xfrac}
\usepackage[unicode,hidelinks,bookmarks=false]{hyperref}
\newcommand{\ie}{i.e.\ }
\newcommand{\cf}{cf.\ }
\newcommand{\resp}{respectively\ }
\newcommand{\Subsection}{\subsection}
\providecommand{\nobreakdash}{\nobreak\hbox{-}}
\setlength{\emergencystretch}{2em}
\multlinegap0pt
\usepackage{bm}
\usepackage{bbm}
\usepackage{dsfont}
\usepackage{xspace}
\usepackage{cancel}
\usepackage{mathtools}
\usepackage{amsthm}
\makeatletter
\@ifundefined{theorem}{\newtheorem{theorem}{Theorem}}{}
\@ifundefined{proposition}{\newtheorem{proposition}[theorem]{Proposition}}{}
\makeatother
\title[Infinitely many solutions for a boundary Yamabe problem]{Infinitely many solutions for a boundary Yamabe problem}
\author{Luca Battaglia, Giusi Vaira, Yixing Pu}
\date{Source: arXiv:2503.06192v2 (CC BY 4.0)}
\begin{document}
\maketitle

\noindent\emph{Source and licence.} Infinitely many solutions for a boundary Yamabe problem. Version arXiv:2503.06192v2 (CC BY 4.0), distributed under the Creative Commons Attribution 4.0 International licence, as recorded in the arXiv licence metadata for this exact version and shown on the abstract page https://arxiv.org/abs/2503.06192v2. Full rights notice: licenses/arxiv-2503.06192-CC-BY-4.0.txt.\par\smallskip

\noindent\textit{Editorial scope.} Counted unit: the Proposition whose source label is \texttt{None} in the article (source0.tex lines 451--453, source label \texttt{None}), delivered together with the article's own complete original proof of it (source0.tex lines 454--506, one \texttt{proof} environment of 53 lines). No mathematics is written, repaired or re-derived: statement and proof are the article's own text line for line. Required context retained uncounted: eq.Z\_i\_1\_2\_expansion\_1 (lines 243--246); eq.Z\_i\_1\_2\_expansion\_2 (lines 243--246). Every definition, display and label of those lines is retained unchanged. Omitted from this package and not redistributed: 1--242, 247--450, 507--934. Nothing inside a retained excerpt is omitted or altered: every retained body is delivered exactly as the article's lines are written, so no omission receipt of any kind accompanies this package. Citations: the retained excerpts carry no citation command at all, so no bibliography entry of the article is transcribed and no reference is redistributed. Reference adaptation: none was needed and none was made; every reference inside the delivered unit resolves inside it, so no unresolved reference or citation is produced. Printed slips kept literal: no printed word, symbol, hypothesis or number of the counted statement or of its proof was altered. Layout and class: the article's own class is \texttt{amsart} with packages including amsmath, amsfonts, amsthm, amsopn, color, amssymb, enumitem, geometry, palatino, graphicx, hyperref, cite, relsize, esint; the wrapper is the lane's pinned \texttt{amsart} editorial wrapper at 10pt on A4, which restates the theorem-like environments the retained text needs and adds the article's own macro definitions that the retained text needs (none), with extra wrapper packages (bm, bbm, dsfont, xspace, cancel, mathtools). The delivered numbering is the unit's own. Assets: no retained span contains a figure environment, a graphics-inclusion command, a PDF-page inclusion, a table or a TikZ picture; no image file is copied into this package, the package declares no asset dependency and no image file is redistributed with it. Licence: version arXiv:2503.06192v2 of this article is distributed under the Creative Commons Attribution 4.0 International licence, as recorded in the arXiv licence metadata for this exact version and shown on the abstract page https://arxiv.org/abs/2503.06192v2. A full rights notice is delivered at licenses/arxiv-2503.06192-CC-BY-4.0.txt. Characters: every retained excerpt is pure ASCII, so the delivered engine, XeLaTeX with its default Latin Modern text font, reports no missing character.
\begin{align}\label{eq.Z_i_1_2_expansion_1}
  &Z_{i,1} := \frac{\partial U_{\mathbf{x}_i,\Lambda}}{\partial r} = U_{\mathbf{x}_i,\Lambda}\frac{(N-2)\Lambda^2 \left(y-\mathbf{x}_i\right)}{{\Lambda^2|\bar{y}-\bar{\mathbf{x}}_{i}|}^2+{(\Lambda y_N+ \mathfrak{D})}^2-1} \cdot \frac{\mathbf{x}_i}{r},\\\nonumber\\ \label{eq.Z_i_1_2_expansion_2}
  &Z_{i,2} := \frac{\partial U_{\mathbf{x}_i,\Lambda}}{\partial \Lambda} =U_{\mathbf{x}_i,\Lambda} \frac{N-2}{2\Lambda} \cdot \frac{\mathfrak{D}^2-\Lambda^2 y_N^2-\Lambda^2\left|\bar{y}-\bar{\mathbf{x}}_i\right|^2-1}{{\Lambda^2|\bar{y}-\bar{\mathbf{x}}_{i}|}^2+{(\Lambda y_N+ \mathfrak{D})}^2-1}.
\end{align}

\begin{proposition}
If $(r, \Lambda)$ is a critical point of the reduced functional $F(r, \Lambda)$ then $\mathfrak c_1=\mathfrak c_2=0$.
\end{proposition}

\begin{proof}
Since $\frac{\partial F}{\partial r}= \frac{\partial F}{\partial \Lambda}=0 \vspace{5pt}$, we only need to show the non-degeneracy of the coefficient matrix with respect to $(\mathfrak c_1,\mathfrak c_2)$ in the following system.
  \begin{align*}
    \begin{pmatrix}
      M_{1,1} & M_{1,2}\\
      M_{2,1} & M_{2,2}
    \end{pmatrix}
    \begin{pmatrix}
    \mathfrak c_1 \vspace{5pt}\\ \mathfrak c_2
    \end{pmatrix} =
    \begin{pmatrix}
      \frac{\partial F}{\partial r} \vspace{10pt}\\
      \frac{\partial F}{\partial \Lambda}
    \end{pmatrix},
  \end{align*}
  where, for $\ell=1,2$,
  \begin{align*}
    \begin{aligned}
      &M_{1,\ell} := \int_{\mathbb{R}_{+}^{N}}\sum_{i=1}^{k} U_{\mathbf{x},\Lambda}^{\frac{4}{N-2}} Z_{i,\ell}\left(\frac{\partial W_{r,\Lambda}}{\partial r}+\frac{\partial \phi}{\partial r}\right)+\int_{\partial\mathbb{R}_{+}^{N}}\sum_{i=1}^{k} U_{\mathbf{x},\Lambda}^{\frac{2}{N-2}} Z_{i,\ell}\left(\frac{\partial W_{r,\Lambda}}{\partial r}+\frac{\partial \phi}{\partial r}\right),\\
      &M_{2,\ell} :=  \int_{\mathbb{R}_{+}^{N}}\sum_{i=1}^{k} U_{\mathbf{x},\Lambda}^{\frac{4}{N-2}} Z_{i,\ell}\left(\frac{\partial W_{r,\Lambda}}{\partial \Lambda}+\frac{\partial \phi}{\partial \Lambda}\right)+\int_{\partial\mathbb{R}_{+}^{N}}\sum_{i=1}^{k} U_{\mathbf{x},\Lambda}^{\frac{2}{N-2}} Z_{i,\ell}\left(\frac{\partial W_{r,\Lambda}}{\partial \Lambda}+\frac{\partial \phi}{\partial \Lambda}\right).
    \end{aligned}
  \end{align*}
  From \eqref{eq.Z_i_1_2_expansion_1} and \eqref{eq.Z_i_1_2_expansion_2}, by direct calculation,  we deduce that
  \begin{align*}
    \left|\frac{\partial Z_{i,\ell}}{\partial r}\right| \leq C U_{\mathbf{x}_{i},\Lambda} \quad\text{and}\quad \left|\frac{\partial Z_{i,\ell}}{\partial \Lambda}\right| \leq C U_{\mathbf{x}_{i},\Lambda}, \quad\text{for}~ \ell=1,2.
  \end{align*}
  Then
$$
  \begin{aligned}
&\int_{\mathbb{R}_{+}^{N}}\sum_{i=1}^{k} U_{\mathbf{x},\Lambda}^{\frac{4}{N-2}} Z_{i,1}\left(\frac{\partial W_{r,\Lambda}}{\partial r}+\frac{\partial \phi}{\partial r}\right)=  \int_{\mathbb{R}_{+}^{N}}\sum_{i=1}^{k} U_{\mathbf{x},\Lambda}^{\frac{4}{N-2}} Z_{i,1}\sum_{j=1}^{k}Z_{j,1}+ \int_{\mathbb{R}_{+}^{N}}\sum_{i=1}^{k} U_{\mathbf{x},\Lambda}^{\frac{4}{N-2}} Z_{i,1}\frac{\partial \phi}{\partial r}\\
  &\quad = \int_{\mathbb{R}_{+}^{N}}\sum_{i=1}^{k} U_{\mathbf{x},\Lambda}^{\frac{4}{N-2}} Z_{i,1}Z_{1,1}+\mathcal O\Biggl(\sum_{i=1}^{k}\sum_{j\neq i} \frac{1}{|\mathbf{x}_i-\mathbf{x}_{j}|^{N-2}}\Biggr) + \int_{\mathbb{R}_{+}^{N}}\sum_{i=1}^{k}\phi \frac{\partial }{\partial r}\left(U_{\mathbf{x},\Lambda}^{\frac{4}{N-2}} Z_{i,1}\right)\\
  &\quad = k \int_{\mathbb{R}_{+}^{N}} U_{\mathbf{x},\Lambda}^{\frac{4}{N-2}} Z_{1,1}^2 + k \mathcal O\left(\frac{1}{\mu}\right)^{\mathfrak m} + O\left(\|\phi\|_{*}\int_{\mathbb{R}_{+}^{N}}\sum_{i=1}^{k}\frac{1}{(1+|y-\mathbf{x}_{i}|)^{N}}\sum_{j}^{k} \frac{1}{(1+|y-\mathbf{x}_{j}|)^{\frac{N+2}{2}+\tau}} d y\right)\\
  &\quad =  k \int_{\mathbb{R}_{+}^{N}} U_{\mathbf{x},\Lambda}^{\frac{4}{N-2}} Z_{1,1}^2+ k \mathcal O\left(\frac{1}{\mu}\right)^{\mathfrak m} + kO\left(\frac{1}{\mu}\right)^{\frac{\mathfrak m}{2}+\sigma} \\
  &\quad = k \int_{\mathbb{R}_{+}^{N}} U_{\mathbf{x},\Lambda}^{\frac{4}{N-2}} Z_{1,1}^2 + k\mathcal O\left(\frac{1}{\mu}\right)^{\frac{\mathfrak m}{2}+\sigma}.
  \end{aligned}
$$
Similarly,
\begin{align*}
  \begin{aligned}
    &\int_{\partial\mathbb{R}_{+}^{N}}\sum_{i=1}^{k} U_{\mathbf{x},\Lambda}^{\frac{2}{N-2}} Z_{i,1}\left(\frac{\partial W_{r,\Lambda}}{\partial r}+\frac{\partial \phi}{\partial r}\right)=   k \int_{\partial\mathbb{R}_{+}^{N}} U_{\mathbf{x},\Lambda}^{\frac{2}{N-2}} Z_{1,1}^2 + k\mathcal O\left(\frac{1}{\mu}\right)^{\frac{\mathfrak m}{2}+\sigma}.
\end{aligned}
\end{align*}
Thus, 
\[M_{1,1}= k \int_{\mathbb{R}_{+}^{N}} U_{\mathbf{x},\Lambda}^{\frac{4}{N-2}} Z_{1,1}^2 + k \int_{\partial\mathbb{R}_{+}^{N}} U_{\mathbf{x},\Lambda}^{\frac{2}{N-2}} Z_{1,1}^2 + k\mathcal O\left(\frac{1}{\mu}\right)^{\frac{\mathfrak m}{2}+\sigma}.\]
Using the same discussion, we have
\begin{align*}
  \begin{aligned}
    &M_{2,2}= k \int_{\mathbb{R}_{+}^{N}} U_{\mathbf{x},\Lambda}^{\frac{4}{N-2}} Z_{1,2}^2 + k \int_{\partial\mathbb{R}_{+}^{N}} U_{\mathbf{x},\Lambda}^{\frac{2}{N-2}} Z_{1,2}^2 + k\mathcal O\left(\frac{1}{\mu}\right)^{\frac{\mathfrak m}{2}+\sigma},\\
    &M_{1,2} = k\mathcal O\left(\frac{1}{\mu}\right)^{\frac{\mathfrak m}{2}+\sigma}, \quad  M_{2,1} = k\mathcal O\left(\frac{1}{\mu}\right)^{\frac{\mathfrak m}{2}+\sigma}.
  \end{aligned}
\end{align*}
This completes the proof.%By \eqref{eq.c_coefficient_matrix_1}, \eqref{eq.c_coefficient_matrix_2} and \eqref{eq.c_coefficient_matrix_3}, we have, for $\ell=1,2$, \begin{align*} \begin{aligned} &\int_{\mathbb{R}_{+}^{n}}\sum_{i=1}^{k} U_{\mathbf{x},\Lambda}^{\frac{4}{n-2}} Z_{i,\ell}\frac{\partial \phi}{\partial \Lambda}\\ &= -\int_{\mathbb{R}_{+}^{n}}\sum_{i=1}^{k}\phi \frac{\partial }{\partial \Lambda}\left(U_{\mathbf{x},\Lambda}^{\frac{4}{n-2}} Z_{i,\ell}\right)\\ &= O\left(\|\phi\|_{*}\int_{\mathbb{R}_{+}^{n}}\sum_{i=1}^{k}\frac{1}{(1+|y-\mathbf{x}_{i}|)^{n}}\sum_{j}^{k} \frac{1}{(1+|y-\mathbf{x}_{j}|)^{\frac{n+2}{2}+\tau}} d y \right)\\ & =kO\Bigl(\frac{1}{\mu}\Bigr)^{\frac{m}{2}+\sigma},\end{aligned}\end{align*} Similarly, for $\ell=1,2$, it holds \begin{align*}\begin{aligned}&\int_{\partial\mathbb{R}_{+}^{n}}\sum_{i=1}^{k} U_{\mathbf{x},\Lambda}^{\frac{2}{n-2}} Z_{i,\ell}\frac{\partial \phi}{\partial r} = kO\Bigl(\frac{1}{\mu}\Bigr)^{\frac{m}{2}+\sigma}.\end{aligned}\end{align*}
\end{proof}

\end{document}
